\documentclass[11pt]{article}
\usepackage[a4paper,margin=18mm]{geometry}
\usepackage{fontspec}
\setmainfont{Latin Modern Roman}
\usepackage{amsmath,amssymb,amsthm,amscd,graphicx,color}
\usepackage{xurl}
\usepackage[unicode,hidelinks]{hyperref}
\allowdisplaybreaks
\setlength\emergencystretch{3em}
\numberwithin{equation}{section}

\newtheorem{Theorem}{Theorem}[section]
\newtheorem{Corollary}[Theorem]{Corollary}
\newtheorem{Lemma}[Theorem]{Lemma}
\newtheorem{Proposition}[Theorem]{Proposition}
{\theoremstyle{definition}
\newtheorem{Note}[Theorem]{Note}
\newtheorem{Definition}[Theorem]{Definition}
\newtheorem{Problem}[Theorem]{Problem}
}

\begin{document}
\begin{center}\Large The center of the universal Askey-Wilson algebra over a field is freely spanned by monomials in its Casimir and three specified central generators when the nonzero deformation parameter is not a root of unity\end{center}
\noindent\textbf{Stable ID:} 2574\par
\noindent\textbf{Authors:} Paul Terwilliger\par
\noindent\textbf{Original work:} The Universal Askey-Wilson Algebra\par
\noindent\textbf{Version:} Official SIGMA 7 (2011) article 069, DOI 10.3842/SIGMA.2011.069; journal PDF and native TeX acquired 2026-10-01, exact bytes pinned by SHA256\par
\noindent\textbf{Selected task scope:} Theorem8.2: F is a field,q is a nonzero element of F not a root of unity,and Delta is exactly the universal Askey-Wilson algebra of Definition1.2. With alpha,beta,gamma the normalized central elements of Definition1.3 and Omega the Casimir of Lemma6.1, the monomials Omega\textasciicircum{}ell alpha\textasciicircum{}r beta\textasciicircum{}s gamma\textasciicircum{}t for all nonnegative ell,r,s,t form an F-vector-space basis of Z(Delta). Algebraic independence and polynomial-algebra description are consequences of this same selected classification, not separate counts.\par
\noindent\textbf{Difficulty:} H2: The generators have noncommuting quadratic relations, so centrality cannot be read off from a commutative presentation. The author establishes a confluent ordered monomial basis, changes its filtered leading terms to include the Casimir, and analyzes commutators of a hypothetical extra central element. A minimal-degree argument then removes every possible leading term, yielding a complete center classification rather than merely exhibiting known central elements.\par
\noindent\textbf{Editorial boundary note (not author text):} The full original Theorem8.2 is retained over the authored field with nonzero q not a root of unity. The full universal relations and central generators, cyclic automorphism, Diamond-Lemma ambiguity check, filtration and coefficient form, Casimir equality/centrality, replacement basis and leading-commutator/minimal-degree classification are bundled as uncounted core. Bergman Diamond Lemma remains a standard cited prerequisite with its local ambiguity calculation supplied. Algebraic independence and polynomial-center description are consequences of this same unit; no Casimir/basis/helper is separately selected.\par
\noindent\textbf{Source:} \url{https://sigma-journal.com/2011/069/}\quad DOI: \url{https://doi.org/10.3842/SIGMA.2011.069}\par
\noindent\textbf{License and attribution:} Copyright retained by the named authors. Published by SIGMA under CC BY 4.0: \url{https://creativecommons.org/licenses/by/4.0/}. Source: \url{https://sigma-journal.com/about.html}. Adaptation consists of standalone excerpt formatting, cover and numbering restoration; original author language and mathematical text are retained.\par
\noindent\textbf{Editorial symbol notice (not author text):} The algebra is universal with three unspecialized central generators, rather than a numerical-parameter quotient AW(3). The not-root-of-unity hypothesis is essential to the leading-commutator argument and preserved in title and scope.\par
\medskip\hrule\medskip\noindent\textbf{Author text begins below.}\par
\setcounter{section}{1}
\setcounter{subsection}{0}
\setcounter{subsubsection}{0}
\setcounter{Theorem}{0}
\setcounter{equation}{0}
% Original source lines 240--330
Our conventions for the paper are as follows.
An algebra is meant to be associative and have a 1.
A subalgebra has the same 1 as the parent algebra.
We f\/ix a f\/ield $\mathbb F$
and  a nonzero $q \in \mathbb F$ such that $q^4\not=1$.

\begin{Definition}[\protect{\cite[equation~(6.1)]{koro}}]
\label{def:aw3}
Let $a$, $b$, $c$ denote scalars in $\mathbb F$.
Def\/ine the $\mathbb F$-algebra ${\rm AW}={\rm AW}_q(a,b,c)$
 by generators and relations in the following way.
The generators are~$A$,~$B$,~$C$.
The relations assert that
\begin{gather*}
 A+ \frac{qBC-q^{-1}CB}{q^2-q^{-2}},
\qquad B+
\frac{qCA-q^{-1}AC}{q^2-q^{-2}},
\qquad
C+
\frac{qAB-q^{-1}BA}{q^2-q^{-2}}
%\label{eq:structureconst}
\end{gather*}
are equal to
$
a/(q+q^{-1})$,
$b/(q+q^{-1})$,
$c /(q+q^{-1})$
respectively. We call ${\rm AW}$ the
{\it Askey--Wilson algebra} that corresponds to
$a$, $b$, $c$.
\end{Definition}

 We now introduce the algebra $\Delta$.

\begin{Definition}
\label{def:uaw}
Def\/ine an $\mathbb F$-algebra
$\Delta=\Delta_q$
 by generators and relations in the following way.
The generators are~$A$,~$B$,~$C$.
The relations assert that each of
\begin{gather}
 A+ \frac{qBC-q^{-1}CB}{q^2-q^{-2}},
%\label{eq:u1}
\qquad B+
\frac{qCA-q^{-1}AC}{q^2-q^{-2}},
%\label{eq:u2}
\qquad
C+
\frac{qAB-q^{-1}BA}{q^2-q^{-2}}
%\label{eq:u3}
\label{eq:comlist}
\end{gather}
is central in $\Delta$.
We call $\Delta$ the {\it universal Askey--Wilson algebra}.
\end{Definition}





\begin{Definition}
\label{def:abc}
\rm
For the
three central elements in
(\ref{eq:comlist}), multiply each by $q+q^{-1}$ to get~$\alpha$,~$\beta$,~$\gamma$. Thus
\begin{gather}
A+ \frac{qBC-q^{-1}CB}{q^2-q^{-2}}
=
\frac{\alpha}{q+q^{-1}},
\label{eq:u1}
\\
B+
\frac{qCA-q^{-1}AC}{q^2-q^{-2}}
=
\frac{\beta}{q+q^{-1}},
\label{eq:u2}
\\
C+
\frac{qAB-q^{-1}BA}{q^2-q^{-2}}
=
\frac{\gamma}{q+q^{-1}}.
\label{eq:u3}
\end{gather}
Note that each of $\alpha$, $\beta$, $\gamma$ is central in
$\Delta$.
(The purpose of the factor $q+q^{-1}$ is
to make the upcoming formula~(\ref{eq:caspreview})
more attractive.)
\end{Definition}
\setcounter{section}{1}
\setcounter{subsection}{0}
\setcounter{subsubsection}{0}
\setcounter{Theorem}{4}
\setcounter{equation}{4}
% Original source lines 393--397
\begin{gather}
\label{eq:caspreview}
\Omega =
qABC+q^2A^2+q^{-2}B^2+q^2C^2-qA\alpha-q^{-1}B\beta -q C\gamma.
\end{gather}
\setcounter{section}{1}
\setcounter{subsection}{0}
\setcounter{subsubsection}{0}
\setcounter{Theorem}{4}
\setcounter{equation}{5}
% Original source lines 421--495
\section[Another presentation of $\Delta$]{Another presentation of $\boldsymbol{\Delta}$}


A bit later in the paper
we will discuss automorphisms of~$\Delta $.
To facilitate this discussion
we give
another presentation for~$\Delta $ by generators and relations.

\begin{Lemma}
\label{lem:abgpres}
The $\mathbb F$-algebra $\Delta $ is generated by
$A$, $B$, $\gamma$. Moreover
\begin{gather}
C  =  \frac{\gamma}{q+q^{-1}}
-\frac{qAB-q^{-1}BA}{q^2-q^{-2}},
\label{eq:getc}
\\
\alpha  =  \frac{B^2A-(q^2+q^{-2})BAB+AB^2+
(q^2-q^{-2})^2A+(q-q^{-1})^2B\gamma}{(q-q^{-1})(q^2-q^{-2})},
\label{eq:getalpha}
\\
\beta  =  \frac{A^2B-(q^2+q^{-2})ABA+BA^2+
(q^2-q^{-2})^2B+(q-q^{-1})^2A\gamma}{(q-q^{-1})(q^2-q^{-2})}.
\label{eq:getbeta}
\end{gather}
\end{Lemma}

\begin{proof}\sloppy
Line~(\ref{eq:getc})
is from (\ref{eq:u3}).
To get
(\ref{eq:getalpha}),
(\ref{eq:getbeta})
eliminate $C$ in
(\ref{eq:u1}),
(\ref{eq:u2})
using
(\ref{eq:getc}).
\end{proof}

 Recall the notation
\begin{gather*}
\lbrack n \rbrack_q= \frac{q^n-q^{-n}}{q-q^{-1}}, \qquad
n=0,1,2,\ldots.
\end{gather*}




{\samepage

\begin{Theorem}
\label{prop:altpres}
The $\mathbb F$-algebra $\Delta $ has a  presentation by generators
$A$, $B$, $\gamma$ and
relations
 \begin{gather*}
 A^3B-\lbrack 3\rbrack_q A^2BA+\lbrack 3\rbrack_q ABA^2-BA^3
= -\big(q^2-q^{-2}\big)^2(AB-BA),
\\
 B^3A-\lbrack 3\rbrack_q B^2AB+\lbrack 3\rbrack_q BAB^2-AB^3
= -\big(q^2-q^{-2}\big)^2(BA-AB),
\\
 A^2B^2-B^2A^2+\big(q^2+q^{-2}\big)(BABA-ABAB) = -\big(q-q^{-1}\big)^2(AB-BA)\gamma,
\\
  \qquad \gamma A = A \gamma,   \qquad \gamma B=B \gamma.
 \end{gather*}
\end{Theorem}

\begin{proof}
Use Lemma~\ref{lem:abgpres} to
express the def\/ining relations
for $\Delta $ in terms of $A$, $B$, $\gamma$.
\end{proof}}
\setcounter{section}{2}
\setcounter{subsection}{0}
\setcounter{subsubsection}{0}
\setcounter{Theorem}{3}
\setcounter{equation}{3}
% Original source lines 521--597
\section[An action of ${\rm  {PSL}}_2(\mathbb Z)$ on $\Delta$]{An action of $\boldsymbol{{\rm  {PSL}}_2(\mathbb Z)}$ on $\boldsymbol{\Delta}$}

We now consider
some automorphisms of~$\Delta $.
Recall that the modular group
${\rm  {PSL}}_2(\mathbb Z)$ has a~presentation
by generators~$\rho$, $\sigma$ and
relations $\rho^3=1$, $\sigma^2=1$.
See for example \cite{alpern}.


\begin{Theorem}
\label{thm:try}
The group
${\rm  {PSL}}_2(\mathbb Z)$ acts on
$\Delta $ as a group of automorphisms in the following way:
\[
\begin{array}{@{}c| ccc | c c c}
u &  A & B  & C
&  \alpha & \beta  & \gamma
\\
\hline
\rho(u) &  B & C & A
&  \beta  & \gamma & \alpha
\\
\sigma(u) &  B & A & C+\frac{AB-BA}{q-q^{-1}}
&  \beta & \alpha & \gamma
\end{array}
\]
\end{Theorem}

\begin{proof}
By Def\/inition
\ref{def:uaw} there exists an automorphism
$P$ of $\Delta $ that sends
\begin{gather*}
A \mapsto B,
\qquad
B \mapsto C,
\qquad
C \mapsto A.
\end{gather*}
Observe $P^3=1$. By
(\ref{eq:u1})--(\ref{eq:u3})
the map $P$ sends
\begin{gather*}
\alpha \mapsto \beta,
\qquad
\beta \mapsto \gamma,
\qquad
\gamma \mapsto \alpha.
\end{gather*}
By Theorem
\ref{prop:altpres} there exists
an automorphism
$S$ of $\Delta $ that sends
\begin{gather*}
A \mapsto B,
 \qquad
B \mapsto A,
\qquad
\gamma \mapsto \gamma.
\end{gather*}
Observe $S^2=1$. By
Lemma
\ref{lem:abgpres}
the map $S$
sends
\begin{gather*}
\alpha \mapsto \beta,
\qquad
\beta \mapsto \alpha,
\qquad
C \mapsto C + \frac{AB-BA}{q-q^{-1}}.
\end{gather*}
The result follows.
\end{proof}
\setcounter{section}{3}
\setcounter{subsection}{0}
\setcounter{subsubsection}{0}
\setcounter{Theorem}{13}
\setcounter{equation}{4}
% Original source lines 1061--1572
\section[A basis for $\Delta $]{A basis for $\boldsymbol{\Delta}$}

 In this section we display a basis for the $\mathbb F$-vector
space $\Delta $.

\begin{Theorem}
\label{thm:basis}
The following is a basis for the $\mathbb F$-vector
space $ \Delta $.
\begin{gather}
\label{eq:basis}
A^iB^jC^k \alpha^r\beta^s\gamma^t, \qquad   i,j,k,r,s,t \geq 0.
\end{gather}
\end{Theorem}

\begin{proof}
We invoke  Bergman's Diamond Lemma
\cite[Theorem~1.2]{berg}.
Consider the symbols
\begin{gather}
\label{eq:letters}
A, \quad B, \quad C, \quad \alpha,\quad \beta, \quad \gamma.
\end{gather}
For an integer $n \geq 0$, by a {\it $\Delta $-word of length $n$}
we mean a sequence
 $x_1x_2\cdots x_n$ such that~$x_i$ is listed in
(\ref{eq:letters}) for $1 \leq i \leq n$.
We interpret the
$\Delta $-word
of length zero to be the multiplicative
identity in $\Delta $.
Consider a
 $\Delta $-word
$w=x_1x_2\cdots x_n$.
By an {\it inversion} for $w$
we mean an ordered pair of integers $(i,j)$
such that
$1\leq i<j\leq n$ and
$x_i$ is strictly to the right of
$x_j$ in the list~(\ref{eq:letters}).
For example $CABA$ has 4 inversions
and
$CB^2A$ has 5 inversions.
A $\Delta $-word is called {\it reducible} whenever it has at least one
inversion, and {\it irreducible} otherwise.
The list
(\ref{eq:basis}) consists of the irreducible $\Delta $-words.
For each integer $n\geq 0$ let
$W_n$ denote the set of $\Delta $-words that have
length~$n$.
Let $W=\cup_{n=0}^\infty W_n$ denote the
set of all
$\Delta $-words.
We now def\/ine a partial
order~$<$ on~$W$.
The def\/inition has two aspects.
$(i)$~For all integers  $n> m\geq 0$,
every word in~$W_m$ is less than
every word in~$W_n$, with respect to $<$.
$(ii)$~For an integer $n\geq 0$
the restriction of $<$ to $W_n$ is
described as follows.
Pick $w, w' \in
 W_n$ and write
 $w=x_1x_2\cdots x_n$.
We say that $w$ {\it covers} $w'$ whenever
there exists an integer~$j$ $(2 \leq j \leq n)$
such that $(j-1,j)$ is an inversion for~$w$,
and
$w'$ is obtained from $w$ by interchanging
$x_{j-1},x_j$.
In this case
$w'$ has one fewer inversions than $w$.
Therefore the transitive closure of the covering
relation on $W_n$ is a partial order on $W_n$, and this is
the restriction of
 $<$ to $W_n$.
We have now def\/ined a partial order $<$ on $W$.
By construction this partial order is
a semi-group partial order
\cite[p.~181]{berg} and satisf\/ies the
descending chain condition
\cite[p.~179]{berg}.
By Def\/inition~\ref{def:uaw} and
Def\/inition~\ref{def:abc}
the def\/ining relations for $\Delta $ can be expressed
as follows:
\begin{gather*}
 BA  =  q^2AB+q\big(q^2-q^{-2}\big)C-q\big(q-q^{-1}\big)\gamma,
\\
 CB  =  q^2BC+q\big(q^2-q^{-2}\big)A-q\big(q-q^{-1}\big)\alpha,
\\
 CA  =  q^{-2}AC+q^{-1}\big(q^{-2}-q^2\big)B-q^{-1}\big(q^{-1}-q\big)\beta,
\\
 \alpha A = A \alpha,
  \qquad
\alpha B = B \alpha,
  \qquad
\alpha C = C \alpha,
\\
 \beta A = A \beta,
  \qquad
\beta B = B \beta,
  \qquad
\beta C = C \beta,
\\
 \gamma A = A \gamma,
  \qquad
\gamma B = B \gamma,
  \qquad
\gamma C = C \gamma,
\\
 \beta \alpha = \alpha \beta,
  \qquad
\gamma \beta = \beta \gamma,
  \qquad
\gamma \alpha = \alpha \gamma.
\end{gather*}
The above equations give reduction rules
for $\Delta$-words,
as we now explain.
Let $w=x_1x_2\cdots x_n$ denote a  reducible $\Delta $-word.
Then there exists an integer $j$ $(2 \leq j \leq n)$
such that $(j-1,j)$ is an inversion for~$w$.
In the above list of equations,
there exists an equation
with $x_{j-1}x_j$ on the left-hand side;
in~$w$ we eliminate $x_{j-1}x_j$ using this equation
and thereby express $w$ as a linear combination of
$\Delta $-words, each less than $w$ with respect to $<$.
Therefore
the reduction rules are compatible
with $<$ in the sense of Bergman \cite[p.~181]{berg}.
In order to employ the Diamond Lemma,
we must show that
the ambiguities
are resolvable in the sense of Bergman
\cite[p.~181]{berg}.
There are potentially two kinds of ambiguities; inclusion
ambiguities and overlap ambiguities
\cite[p.~181]{berg}.
For the present example there are no inclusion ambiguities.
The only nontrivial overlap ambiguity
involves the word
$CBA$. This word can be reduced
in two ways; we could evaluate $CB$ f\/irst or we could
evaluate $BA$ f\/irst.
 Either way, after a three-step reduction
we obtain the same result, which is
\begin{gather*}
 q^{-1}CBA = qABC
+\big(q^2-q^{-2}\big)A^2
-\big(q^2-q^{-2}\big)B^2
+\big(q^2-q^{-2}\big)C^2
\\
\phantom{q^{-1}CBA =}{}
- \big(q-q^{-1}\big)A \alpha
+ \big(q-q^{-1}\big)B \beta
- \big(q-q^{-1}\big)C \gamma .
\end{gather*}
Therefore the
overlap ambiguity $CBA$ is re\-sol\-vable.
We conclude that every  ambiguity is
re\-sol\-vable, so
by the Diamond Lemma
\cite[Theorem~1.2]{berg}
the elements
(\ref{eq:basis}) form a basis for
$\Delta $.
\end{proof}

 On occasion we
wish to discuss the coef\/f\/icients
when an element of $\Delta$ is written as a linear
combination of the elements
(\ref{eq:basis}). To facilitate this discussion
we def\/ine a bilinear form
$(\,,) : \Delta \times \Delta \to \mathbb F$
such that $( u,v) = \delta_{u,v}$
for all elements $u$, $v$ in the basis
(\ref{eq:basis}). In other words
the basis
(\ref{eq:basis}) is
orthonormal with respect to
$( \,,) $.
 Observe that
$(\,,)$ is symmetric.
For $u \in \Delta$,
\begin{gather}
u = \sum \big(u, A^iB^jC^k\alpha^r\beta^s\gamma^t \big)
 A^iB^jC^k\alpha^r\beta^s\gamma^t,
\label{eq:bilsum}
\end{gather}
where the sum is over all elements
 $A^iB^jC^k\alpha^r\beta^s\gamma^t$ in the basis
(\ref{eq:basis}).

\begin{Definition}
\label{def:contrib}
\rm
Let $u \in \Delta $. A given element
 $A^iB^jC^k\alpha^r\beta^s\gamma^t$ in the basis
(\ref{eq:basis}) is said to
{\it contribute to
  $u$} whenever $(u,
 A^iB^jC^k\alpha^r\beta^s\gamma^t )
  \not=0$.
\end{Definition}







\section[A filtration of $\Delta$]{A f\/iltration of $\boldsymbol{\Delta}$}


In this section we obtain
a f\/iltration of $\Delta $
which
is related to the basis from Theorem~\ref{thm:basis}.
 This f\/iltration will be useful
when we investigate the center  $Z(\Delta)$ later in the paper.


We recall some notation.
For subspaces $H$, $K$ of $\Delta $ def\/ine
$HK={\rm Span}\lbrace hk | h \in H, \; k \in K\rbrace$.

\begin{Definition}
\label{def:ui}
We def\/ine subspaces
$\lbrace \Delta_n\rbrace_{n=0}^\infty$
of $\Delta $ such that
\begin{gather*}
\Delta_0=\mathbb F 1,
  \qquad
\Delta_1 =\Delta_0+{\rm Span}\lbrace A,B,C,\alpha,\beta,\gamma\rbrace,
 \qquad
\Delta_n=\Delta_1 \Delta_{n-1}, \qquad n=1,2,\ldots
\end{gather*}
\end{Definition}

\begin{Lemma}
\label{lem:filt}
The following $(i)$--$(iii)$ hold.
\begin{enumerate}\itemsep=0pt
\item[$(i)$]
 $\Delta_{n-1}\subseteq \Delta_n$ for  $n\geq 1$.
\item[$(ii)$]
$\Delta=\cup_{n=0}^\infty \Delta_n$.
\item[$(iii)$]
$\Delta_m \Delta_n = \Delta_{m+n} $ for
 $m,n\geq 0$.
\end{enumerate}
\end{Lemma}

\begin{proof}
$(i)$ Since $\Delta_n=\Delta_1 \Delta_{n-1}$
and $1 \in \Delta_1$.
$(ii)$ Since $A$, $B$, $C$, $\alpha$, $\beta$, $\gamma $ generate~$\Delta$.
$(iii)$ Each side is equal to
$(\Delta_1)^{m+n}$.
\end{proof}

By Lemma~\ref{lem:filt} the sequence
$\lbrace \Delta_n\rbrace_{n=0}^\infty$ is a f\/iltration of~$\Delta$ in the sense of
\cite[p.~202]{carter}.


\begin{Lemma}
\label{lem:almostcom}
Each of the following is contained in $\Delta_1$:
\begin{gather*}
qAB-q^{-1}BA,
\qquad
qBC-q^{-1}CB,
\qquad
qCA-q^{-1}AC.
\end{gather*}
\end{Lemma}

\begin{proof}
Each of the three expressions is a linear combination of
$A$, $B$, $C$, $\alpha$, $\beta$, $\gamma$ and these are contained in~$\Delta_1$.
\end{proof}

\begin{Theorem}\label{lem:unbasis}
For all integers $n\geq 0$ the following is
a basis for the~$\mathbb F$-vector space~$\Delta_n$:
\begin{gather}
\label{eq:basislist}
A^iB^jC^k\alpha^r \beta^s \gamma^t, \qquad
i,j,k,r,s,t\geq 0, \qquad
i+j+k+r+s+t \leq n.
\end{gather}
\end{Theorem}

\begin{proof}
The elements~(\ref{eq:basislist}) are linearly independent by
Theorem
\ref{thm:basis}. We show that the
elements~(\ref{eq:basislist}) span $\Delta_n$.
We will use induction on $n$. Assume $n\geq 2$; otherwise
the result holds by Def\/inition~\ref{def:ui}.
By Def\/inition~\ref{def:ui} $\Delta_n$ is
spanned by the set of elements of the form
$x_1x_2\cdots x_n$
where $x_i \in \lbrace 1, A,B,C,\alpha,\beta,\gamma\rbrace$
for $1 \leq i \leq n$.
Therefore $\Delta_n$ is spanned by the set of elements
of the form
$x_1x_2\cdots x_m$ where $0 \leq m\leq n$ and
$x_i \in \lbrace A,B,C,\alpha,\beta,\gamma\rbrace$
for $1 \leq i \leq m$.
Therefore $\Delta_n$ is spanned by $\Delta_{n-1}$ together with
the set of elements
of the form
$x_1x_2\cdots x_n$ where
$x_i \in \lbrace A,B,C,\alpha,\beta,\gamma\rbrace$
for $1 \leq i \leq n$.
Consider such an element
$x_1x_2\cdots x_n$.
By Lemma~\ref{lem:almostcom} and since
each of $\alpha$, $\beta$, $\gamma$ is central,
we f\/ind that
for $2 \leq j \leq n$,
\begin{gather*}
x_1\cdots  x_{j-1}x_j\cdots  x_n \in
\mathbb F x_1\cdots x_{j}x_{j-1}\cdots  x_n + \Delta_{n-1}.
\end{gather*}
By the above comments
$\Delta_n$ is spanned by $\Delta_{n-1}$ together with the
set
\begin{gather*}
A^iB^jC^k\alpha^r \beta^s \gamma^t, \qquad
i,j,k,r,s,t\geq 0, \qquad
i+j+k+r+s+t = n.
\end{gather*}
By this and induction
$\Delta_n$ is spanned by the elements~(\ref{eq:basislist}).
We have shown that
the elements~(\ref{eq:basislist}) form a basis  for the $\mathbb F$-vector space~$\Delta_n$.
\end{proof}

  Let $V$ denote a vector space over
$\mathbb F$ and let $U$ denote a subspace of~$V$.
By a {\it complement of~$U$ in~$V$} we mean
a subspace $U'$ of $V$ such that
$V=U+U'$ (direct sum).


\begin{Corollary}
\label{lem:com}
For all integers $n\geq 1$ the following is a basis
for a complement of $\Delta_{n-1}$ in $\Delta_n$:
\begin{gather*}
%\label{eq:combas}
A^iB^jC^k\alpha^r \beta^s \gamma^t, \qquad
i,j,k,r,s,t\geq 0, \qquad
i+j+k+r+s+t = n.
\end{gather*}
\end{Corollary}
\begin{proof}
Use Theorem
\ref{lem:unbasis}.
\end{proof}



\section[The Casimir element $\Omega$]{The Casimir element $\boldsymbol{\Omega}$}

We turn our attention to the center $Z(\Delta)$.
In this section we discuss a certain
 element $\Omega \in Z(\Delta)$
 called the Casimir element.
The name is
motivated by
\cite[equation~(1.3)]{Zhidd}.
 In Section 7
we will use~$\Omega$ to describe
$Z(\Delta)$.
We acknowledge that the results of this section are
extensions of
\cite[Lemma~1]{posta2},
\cite[Section~6]{koro},
\cite[equation~(1.3)]{Zhidd}.

\begin{Lemma}
\label{lem:six}
The following elements of $\Delta $ coincide:
\begin{gather*}
qABC + q^2A^2+q^{-2}B^2+q^2C^2-q A\alpha -q^{-1} B\beta -q C\gamma,
\\
qBCA +
q^2A^2 +
q^2B^2+q^{-2}C^2
-q A\alpha
-qB\beta -q^{-1}C\gamma,
\\
qCAB +
q^{-2}A^2+q^2B^2
+q^2 C^2
-q^{-1} A\alpha -q B \beta
-q C\gamma,
\\
q^{-1}CBA + q^{-2}A^2+q^{2}B^2+q^{-2}C^2-q^{-1}A\alpha -q B\beta -
q^{-1} C\gamma,
\\
q^{-1}ACB +
q^{-2}A^2 +
q^{-2}B^2+q^{2}C^2
-q^{-1} A\alpha
-q^{-1}B\beta -q C\gamma,
\\
q^{-1}BAC +
q^{2}A^2+q^{-2}B^2
+q^{-2}C^2
-q A\alpha
-q^{-1} B \beta
-q^{-1} C\gamma.
\end{gather*}
We denote this common element by $\Omega$.
\end{Lemma}

\begin{proof}
Denote the displayed sequence of elements by
$\Omega^+_B$,
$\Omega^+_C$,
$\Omega^+_A$,
$\Omega^-_B$,
$\Omega^-_C$,
$\Omega^-_A$.
The automorphism
$\rho$ cyclically permutes
$\Omega^+_A$,
$\Omega^+_B$,
$\Omega^+_C$
and cyclically permutes
$\Omega^-_A$, $\Omega^-_B$, $\Omega^-_C$.
The element
$\Omega^+_B-\Omega^-_C$
is
equal to $(q-q^{-1})A$ times
\begin{gather}
\label{eq:need0}
\big(q+q^{-1}\big)A + \frac{qBC-q^{-1}CB}{q-q^{-1}} -\alpha.
\end{gather}
The element
(\ref{eq:need0}) is zero by Def\/inition
\ref{def:abc} so
$\Omega^+_B=\Omega^-_C$.
In this equation we apply $\rho$ twice to get
$\Omega^+_C=\Omega^-_A$
and
$\Omega^+_A=\Omega^-_B$.
The element
$\Omega^+_B-\Omega^-_A$ is equal to
\begin{gather}
\label{eq:need02}
\big(q+q^{-1}\big)C + \frac{qAB-q^{-1}BA}{q-q^{-1}} -\gamma
\end{gather}
times $(q-q^{-1})C$.
The element
(\ref{eq:need02}) is zero by Def\/inition
\ref{def:abc} so
$\Omega^+_B=\Omega^-_A$.
Applying $\rho$ twice we get
$\Omega^+_C=\Omega^-_B$
and
$\Omega^+_A=\Omega^-_C$.
By these comments
$\Omega^+_B$,
$\Omega^+_C$,
$\Omega^+_A$,
$\Omega^-_B$,
$\Omega^-_C$,
$\Omega^-_A$ coincide.
\end{proof}

\begin{Theorem}
\label{lem:center}
The element $\Omega$ from
Lemma~{\rm \ref{lem:six}} is central in $\Delta $.
\end{Theorem}

\begin{proof}
We f\/irst show $\Omega A = A \Omega $.
We will work with the equations~(\ref{eq:u2}),
(\ref{eq:u3}) from Def\/inition~\ref{def:abc}. Consider the equation which is
$qC $ times~(\ref{eq:u2}) plus~(\ref{eq:u2}) times $q^{-1}C$
minus $\gamma$ times~(\ref{eq:u2})
plus $\beta$ times~(\ref{eq:u3})
minus
$q^{-1}B$ times~(\ref{eq:u3})
minus~(\ref{eq:u3}) times~$qB$.
After some cancellation this equation yields
$\Omega^+_B A- A\Omega^+_C =0$, where
$\Omega^+_B$, $\Omega^+_C$
are from the proof of Lemma
\ref{lem:six}.
Therefore
$\Omega A= A\Omega $.
 One similarly
f\/inds
$
\Omega B = B \Omega$ and
$
\Omega C = C \Omega$.
The elements $A$, $B$, $C$ generate~$\Delta $ so~$\Omega$ is central in~$\Delta$.
\end{proof}
\setcounter{section}{6}
\setcounter{subsection}{0}
\setcounter{subsubsection}{0}
\setcounter{Theorem}{4}
\setcounter{equation}{2}
% Original source lines 1618--2024
\section[A basis for~$\Delta $ that involves~$\Omega$]{A basis for $\boldsymbol{\Delta}$ that involves~$\boldsymbol{\Omega}$}

 In Theorem~\ref{thm:basis} we displayed a basis for
$\Delta $. In this section we display a related
basis for $\Delta $ that involves the Casimir element~$\Omega$.
In the next section we will use the related basis
to describe  the center~$Z(\Delta )$.

 Recall the f\/iltration $\lbrace \Delta_n \rbrace_{n=0}^\infty$
of $\Delta $ from Def\/inition~\ref{def:ui}.


\begin{Lemma}
\label{lem:omegapower}
For all integers $\ell \geq 1$
the following hold:
\begin{enumerate}\itemsep=0pt
\item[$(i)$] $\Omega^\ell \in \Delta_{3 \ell}$.
\item[$(ii)$] $\Omega^\ell -q^{\ell^2}A^\ell B^\ell C^\ell  \in
\Delta_{3\ell-1}$.
\end{enumerate}
\end{Lemma}

\begin{proof}
Consider the expression for $\Omega$ from the
f\/irst displayed line of
Lemma~\ref{lem:six}.
The term $qABC$ is in $\Delta_3$
and the remaining terms are in $\Delta_2$. Therefore
$\Omega \in \Delta_3$ and $\Omega-qABC \in \Delta_2$. By this and
Lemma~\ref{lem:filt}$(iii)$ we f\/ind
$\Omega^\ell \in \Delta_{3\ell}$ and
$\Omega^\ell - q^\ell (ABC)^\ell
\in \Delta_{3\ell-1}$.
Using Lemma~\ref{lem:almostcom}  we obtain
$(ABC)^\ell -q^{\ell (\ell-1)}A^\ell B^\ell C^\ell \in \Delta_{3\ell -1}$.
By these comments
$\Omega^\ell -q^{\ell^2}A^\ell B^\ell C^\ell  \in \Delta_{3\ell-1}$.
\end{proof}

\begin{Lemma}
\label{lem:basis2c}
For all integers $n\geq 1$
the following is a basis
for a complement of
$\Delta_{n-1}$ in $\Delta_n$:
\begin{gather*}
 A^iB^jC^k\Omega^\ell \alpha^r \beta^s \gamma^t, \qquad
i,j,k,\ell, r,s,t\geq 0,
\qquad ijk=0,
 \qquad
i+j+k+3\ell+r+s+t = n.
\end{gather*}
\end{Lemma}

\begin{proof}
Let $\mathbb I_n$ denote the set consisting of the 6-tuples
$(i,j,k,r,s,t)$ of nonnegative integers whose sum is~$n$.
By Corollary~\ref{lem:com} the
following
is a basis for a complement of
$\Delta_{n-1}$ in $\Delta_n$:
\begin{gather*}
A^iB^jC^k\alpha^r\beta^s\gamma^t, \qquad
(i,j,k,r,s,t) \in {\mathbb I}_n.
\end{gather*}
Let $\mathbb J_n$ denote the set consisting of the 7-tuples
$(i,j,k,\ell,r,s,t)$ of nonnegative integers such that
$ijk=0$ and
$i+j+k+3\ell+r+s+t = n$.
Observe that the map
\begin{gather*}
\mathbb J_n  \to  \mathbb I_n,
\qquad
(i,j,k,\ell,r,s,t)  \mapsto   (i+\ell,j+\ell,k+\ell,r,s,t)
\end{gather*}
is a bijection.
Suppose we are given $(i,j,k,\ell,r,s,t) \in \mathbb J_n$.
By
Lemma \ref{lem:omegapower},
$\Delta_{n-1}$ contains
\begin{gather*}
A^iB^jC^k\Omega^\ell \alpha^r \beta^s \gamma^t -
q^{\ell^2} A^iB^jC^k A^\ell B^\ell C^\ell \alpha^r \beta^s \gamma^t.
\end{gather*}
By Lemma~\ref{lem:almostcom},  $\Delta_{n-1}$ contains
\begin{gather*}
A^iB^jC^k A^\ell B^\ell C^\ell \alpha^r \beta^s \gamma^t
-
q^{2j \ell} A^{i+\ell}B^{j+\ell}C^{k+\ell}  \alpha^r \beta^s \gamma^t.
\end{gather*}
Therefore $\Delta_{n-1}$ contains
\begin{gather*}
A^iB^jC^k\Omega^\ell \alpha^r \beta^s \gamma^t -
q^{\ell(2j+ \ell)} A^{i+\ell}B^{j+\ell}C^{k+\ell}  \alpha^r \beta^s \gamma^t.
\end{gather*}
By these comments the following is a basis for a complement of
$\Delta_{n-1}$ in $\Delta_n$:
\begin{gather*}
A^iB^jC^k\Omega^\ell \alpha^r\beta^s\gamma^t, \qquad
(i,j,k,\ell,r,s,t) \in \mathbb J_n.
\end{gather*}
The result follows.
\end{proof}

\begin{Note}
Pick an integer $n \geq 1$.
In
Corollary
\ref{lem:com} and
Lemma
\ref{lem:basis2c}
we mentioned a
complement of $\Delta_{n-1}$ in $\Delta_n$. These
complements are not the same
in general.
\end{Note}


\begin{Proposition}
\label{lem:basis2n}
For all integers $n\geq 0$
the following is a basis
for
the $\mathbb F$-vector space
$\Delta_n$:
\begin{gather*}
 A^iB^jC^k\Omega^\ell \alpha^r \beta^s \gamma^t, \qquad
i,j,k,\ell, r,s,t\geq 0,
\qquad ijk=0,
 \qquad
i+j+k+3\ell+r+s+t \leq  n.
\end{gather*}
\end{Proposition}

\begin{proof}
By Lemma~\ref{lem:basis2c} and
$\Delta_0=\mathbb F 1$.
\end{proof}




\begin{Theorem}
\label{lem:basis2}
The following is a basis for the $\mathbb F$-vector space
$\Delta$:
\begin{gather*}
A^iB^jC^k\Omega^\ell \alpha^r \beta^s \gamma^t, \qquad
i,j,k,\ell, r,s,t\geq 0,
 \qquad ijk=0.
\end{gather*}
\end{Theorem}

\begin{proof}
Combine
Lemma~\ref{lem:filt}$(ii)$
and
Proposition~\ref{lem:basis2n}.
\end{proof}







\section[The center $Z(\Delta)$]{The center $\boldsymbol{Z(\Delta)}$}

In this section we give a detailed description of
the center $Z(\Delta)$, under the assumption that $q$ is not a root of
 unity.
For such $q$ we show that $Z(\Delta)$ is generated by
$\Omega$, $\alpha$, $\beta$, $\gamma$ and
 isomorphic to a polynomial algebra in four
variables.

  Recall the commutator
 $\lbrack r,s\rbrack=rs-sr$.


\begin{Lemma}
\label{lem:ijk}
Let
$i$, $j$, $k$
denote nonnegative integers.
Then $\Delta_{i+j+k}$ contains each of the fol\-lo\-wing:
\begin{gather}
\label{eq:1}
 \big\lbrack A,A^iB^jC^k\big\rbrack - \big(1-q^{2j-2k}\big)A^{i+1}B^jC^k,
\\
\label{eq:2}
 \big\lbrack B,A^iB^jC^k\big\rbrack - \big(q^{2i}-q^{2k}\big)A^iB^{j+1}C^k,
\\
\label{eq:3}
 \big\lbrack C,A^iB^jC^k\big\rbrack - \big(q^{2j-2i}-1\big)A^iB^jC^{k+1}.
\end{gather}
\end{Lemma}

\begin{proof}
Concerning (\ref{eq:1}), observe
\begin{gather*}
\big\lbrack A,A^iB^jC^k\big\rbrack = A^{i+1}B^jC^k -A^iB^jC^kA.
\end{gather*}
By Lemma
\ref{lem:almostcom} $\Delta_{i+j+k}$ contains
\begin{gather*}
A^iB^jC^kA - q^{2j-2k}A^{i+1}B^jC^k.
\end{gather*}
By these comments
$\Delta_{i+j+k}$ contains
(\ref{eq:1}).
One similarly f\/inds that
$\Delta_{i+j+k}$ contains
(\ref{eq:2}),
(\ref{eq:3}).
\end{proof}


\begin{Theorem}
\label{thm:centerbasis}
Assume that $q$ is not a root of unity. Then the following is a basis for
the $\mathbb F$-vector space~$Z(\Delta)$.
\begin{gather}
\Omega^\ell \alpha^r \beta^s\gamma^t, \qquad  \ell,r,s,t \geq 0.
\label{eq:cbasis}
\end{gather}
\end{Theorem}

\begin{proof}
Abbreviate $Z=Z(\Delta)$.
The elements~(\ref{eq:cbasis})
are linearly independent by
Theorem~\ref{lem:basis2},
so it suf\/f\/ices to
show that they span~$Z$.
Let~$Z'$ denote the subspace of~$\Delta$ spanned by~(\ref{eq:cbasis}), and note that
$Z' \subseteq Z$. To show
$Z'= Z$, we assume that $Z'$ is properly contained in
$Z$ and get a contradiction.
Def\/ine the set $E=Z\backslash Z'$
and note $E\not=\varnothing$.
We have
$\Delta_0=\mathbb F 1
\subseteq Z'$ so
$E \cap \Delta_0 =\varnothing$.
By this and
Lemma~\ref{lem:filt}$(i)$,$(ii)$ there exists a unique integer
$n\geq 1$ such that
$E\cap \Delta_n \not=\varnothing$ and
$E\cap \Delta_{n-1} =\varnothing$.
Fix $u \in E\cap \Delta_n$.
Let
$S=S(u)$ denote
the set of 6-tuples $(i,j,k,r,s,t)$ of nonnegative integers
whose sum is $n$ and
 $A^iB^jC^k\alpha^r\beta^s\gamma^t$
contributes to $u$ in
the sense of
Def\/inition~\ref{def:contrib}.
By~(\ref{eq:bilsum}) and Corollary~\ref{lem:com},
$\Delta_{n-1}$ contains
\begin{gather}
\label{eq:dif}
u- \sum_{(i,j,k,r,s,t) \in S}
\big( u,
A^iB^jC^k\alpha^r \beta^s \gamma^t \big)
A^iB^jC^k\alpha^r \beta^s \gamma^t.
\end{gather}
We are going to show that
$i=j=k$ for all
$(i,j,k,r,s,t) \in S$.
To this end
we f\/irst claim that $j=k$ for all
$(i,j,k,r,s,t) \in S$.
To prove the claim,
take the commutator of $A$ with~(\ref{eq:dif}) and evaluate the
result using the following facts.
By construction $u \in E \subseteq Z$
so
  $\lbrack A,u\rbrack=0$.
By Lemma~\ref{lem:filt}$(iii)$
$A\Delta_{n-1} \subseteq  \Delta_n $
and
$\Delta_{n-1}A \subseteq \Delta_n$
so
 $\lbrack A,\Delta_{n-1}\rbrack \subseteq \Delta_n$.
Moreover
 each of~$\alpha$,~$\beta$,~$\gamma$
 is central.
The above evaluation shows that~$\Delta_n$ contains
\begin{gather*}
\sum_{(i,j,k,r,s,t) \in S}
\big( u, A^iB^jC^k
\alpha^r \beta^s \gamma^t \big)
\big\lbrack A, A^iB^jC^k \big\rbrack
\alpha^r \beta^s \gamma^t.
\end{gather*}
Pick
$(i,j,k,r,s,t) \in S$.
By
Lemma~\ref{lem:filt}$(iii)$ and
Lemma~\ref{lem:ijk}, $\Delta_n$ contains
\begin{gather*}
\big\lbrack A, A^iB^jC^k \big\rbrack
\alpha^r \beta^s \gamma^t
-
\big(1-q^{2j-2k}\big)  A^{i+1}B^jC^k
\alpha^r \beta^s \gamma^t.
\end{gather*}
By the above comments
$\Delta_n$ contains
\begin{gather*}
\sum_{
(i,j,k,r,s,t) \in S}
\big(u, A^iB^jC^k
\alpha^r \beta^s \gamma^t \big)
\big(1-q^{2j-2k}\big)  A^{i+1}B^jC^k
\alpha^r \beta^s \gamma^t.
\end{gather*}
For all
$(i,j,k,r,s,t) \in S$
the element
$A^{i+1}B^jC^k
\alpha^r \beta^s \gamma^t$
is contained
in the basis for the complement of $\Delta_n$ in $\Delta_{n+1}$
given in Corollary~\ref{lem:com}. Therefore
\begin{gather*}
\big(u, A^iB^jC^k
\alpha^r \beta^s \gamma^t\big)
\big(1-q^{2j-2k}\big)  = 0
\qquad
\forall \;(i,j,k,r,s,t) \in S.
\end{gather*}
By the def\/inition of $S$ we have
$( u, A^iB^jC^k
\alpha^r \beta^s \gamma^t )\not=0$
 for
all
$(i,j,k,r,s,t) \in S$.
Therefore
$1-q^{2j-2k} = 0 $ for all
$(i,j,k,r,s,t) \in S$.
The scalar $q$ is not a root of unity
so $j=k$ for all
$(i,j,k,r,s,t) \in S$.
The claim is proved.
We next claim that
$i=j$ for all
$(i,j,k,r,s,t) \in S$.
This claim is proved like the previous one,
 except that
as we begin the argument below~(\ref{eq:dif}),
we use~$C$ instead of~$A$ in the commutator.
By the two claims
$i=j=k$ for all
$(i,j,k,r,s,t) \in S$.
In light of this we revisit the assertion above~(\ref{eq:dif}) and conclude that
$\Delta_{n-1}$ contains
\begin{gather*}
u-\sum_{
(i,i,i,r,s,t) \in S}
\big(u, A^iB^iC^i
\alpha^r \beta^s \gamma^t \big)
A^{i}B^iC^i
\alpha^r \beta^s \gamma^t.
\end{gather*}
Pick
$(i,i,i,r,s,t) \in S$.
By
Lemma~\ref{lem:filt}$(iii)$ and
Lemma~\ref{lem:omegapower},
$\Delta_{n-1}$ contains
\begin{gather*}
A^iB^iC^i
\alpha^r \beta^s \gamma^t
-q^{-i^2}\Omega^i \alpha^r \beta^s \gamma^t.
\end{gather*}
By these comments
$\Delta_{n-1}$ contains
\begin{gather*}
\label{eq:summ}
u-
\sum_{
(i,i,i,r,s,t) \in S}
q^{-i^2}
\big(u, A^iB^iC^i
\alpha^r \beta^s \gamma^t  \big)
\Omega^i \alpha^r \beta^s \gamma^t.
\end{gather*}
In the above expression let
$\psi$ denote the main sum,
so that
$u-\psi \in \Delta_{n-1}$.
Observe
$\psi \in Z' \subseteq Z$.
Recall $u \in E=Z\backslash Z'$ so
$u \in Z$ and $u \not\in Z'$. By these comments
$u - \psi \in Z$ and
$u - \psi \not\in Z'$.
Therefore $u - \psi \in E$ so
 $u - \psi \in E\cap \Delta_{n-1}$.
This contradicts
 $E\cap \Delta_{n-1}=\varnothing $ so
$Z=Z'$. The result follows.
\end{proof}
\begin{thebibliography}{99}
\bibitem[1]{alpern}
Alperin R.C.,
Notes: ${\rm  {PSL}}_2(\mathbb Z)= \mathbb Z_2 \star \mathbb Z_3$,
\href{http://dx.doi.org/10.2307/2324963}{{\it Amer. Math. Monthly}} {\bf 100} (1993), 385--386.

\bibitem[16]{berg}
Bergman G.,
The diamond  lemma for ring theory,
\href{http://dx.doi.org/10.1016/0001-8708(78)90010-5}{{\it Adv. Math.}} {\bf 29} (1978), 178--218.


\bibitem[17]{carter}
Carter R.,
Lie algebras of f\/inite and af\/f\/ine type,
\href{http://dx.doi.org/10.1017/CBO9780511614910}{{\it Cambridge Studies in Advanced Mathematics}}, Vol.~96,
Cambridge University Press, Cambridge, 2005.



\bibitem[29]{posta2}
Havl{\'{\i}}{\v{c}}ek M.,  Po\v{s}ta S.,
On the classif\/ication of irreducible f\/inite-dimensional representations of $U'_q({\rm so}_3)$ algebra,
\href{http://dx.doi.org/10.1063/1.1328078}{{\it J. Math. Phys.}} {\bf 42} (2001), 472--500.




\bibitem[40]{koro}
Korovnichenko A., Zhedanov A.,
Classical Leonard triples,
in Elliptic Integrable Systems (Kyoto, 2004), Editors M.~Noumi and K.~Takasaki,
{\it Rokko Lectures in Mathematics}, no.~18, Kobe University, 2005, 71--84.


\bibitem[65]{Zhidd}
Zhedanov A.S.,
``Hidden symmetry'' of the Askey--Wilson polynomials,
\href{http://dx.doi.org/10.1007/BF01015906}{{\it Theoret. and Math. Phys.}} {\bf 89} (1991), 1146--1157.


\end{thebibliography}
\end{document}
